\section{Learning Channel and Activity Denoisers}
\label{sec:learning}

The receiver in Section~\ref{sec:algorithm} requires a channel denoiser
and an activity classifier at each layer. We first learn the channel
score from Gaussian-perturbed channel samples. The pretrained network
then initializes the unfolded layers, which are trained jointly with
the activity classifiers and message calibration parameters.

\subsection{First-Order Score Matching}

Let $h_R\in\mathbb R^{2D}$ denote a normalized channel sample.
We perturb it as $u=h_R+\sqrt c\,\varepsilon$, where $c=v/2$
and $\varepsilon\sim\mathcal N(0,\I_{2D})$. Denote the perturbed
density by $q_{1,R}(u;v)$ and its score by
$s^\star(u,v)=\nabla_u\log q_{1,R}(u;v)$.
Tweedie's formula relates this score to the conditional channel mean
\cite{efron2011,cai2026}:
\begin{equation}
 \E[h_R\mid u]=u+\frac v2\,s^\star(u,v).
 \label{eq:twc-tweedie}
\end{equation}
Equation~\eqref{eq:twc-tweedie} provides the channel estimate from a
score evaluation, avoiding direct computation of the integral in
\eqref{eq:twc-local-targets}.

Since the channel score is unavailable analytically, we learn it
using denoising score matching \cite{vincent2011,song2019}.
The Gaussian perturbation has conditional score
$\nabla_u\log q(u\mid h_R,v)=-\varepsilon/\sqrt{v/2}$.
Substituting this expression into the denoising score matching
objective yields
\begin{equation}
 \mathcal L_{\rm DSM}(\theta)=\E_{h_R,v,\varepsilon}
 \left[\frac{\|\sqrt{v/2}\,s_\theta(h_R+\sqrt{v/2}\varepsilon,v)
                      +\varepsilon\|^2}{2D}\right].
 \label{eq:twc-dsm}
\end{equation}
Training samples are drawn from the channel distribution before
activity masking and perturbed over a range of noise levels.
The resulting network is used to initialize
$s_{\theta_\ell}$ in \eqref{eq:twc-conditional-mean}.

\subsection{Joint Channel and Activity Training}

The two outputs in \eqref{eq:twc-local-targets} have different training
targets: the channel denoiser estimates the channel of an active user,
whereas the classifier estimates its activity probability.
We therefore use a squared channel error on active users and binary
cross-entropy on all users.

Suppose the receiver training set contains $B$ scenes, with $K_i$
candidate users in scene $i$. Let $h_{ik,R}$ and $\alpha_{ik}$
denote the channel and activity label of user $k$ in scene $i$.
Define the total number of users and active-channel energy as
\begin{equation}
 T_{\rm tr}=\sum_{i=1}^{B}K_i,\qquad
 E_{\rm tr}=\sum_{i=1}^{B}\sum_{k=1}^{K_i}
               \alpha_{ik}\|h_{ik,R}\|^2.
 \label{eq:twc-training-statistics}
\end{equation}
The channel and activity losses at layer $\ell$ are respectively
given by
\begin{align}
 R_H^{(\ell)}&=\frac{1}{E_{\rm tr}}\sum_{i,k}
       \alpha_{ik}\|\mu_{ik\ell}-h_{ik,R}\|^2,
 \label{eq:twc-channel-loss}\\
 R_a^{(\ell)}&=-\frac{1}{T_{\rm tr}}\sum_{i,k}
       \left[\alpha_{ik}\log p_{ik\ell}
       +(1-\alpha_{ik})\log(1-p_{ik\ell})\right].
 \label{eq:twc-activity-loss}
\end{align}
The activity mask in \eqref{eq:twc-channel-loss} restricts channel
supervision to active users. The classifier is trained using both
active and inactive users, and its output weights the channel estimate
through \eqref{eq:twc-estimate}. The constants $E_{\rm tr}$ and
$T_{\rm tr}$ normalize the channel energy and user count, respectively.

For a receiver with $m$ layers, the joint training objective is
\begin{equation}
 J_m=\frac{\sum_{\ell=1}^{m}\omega_\ell
                  \left(R_H^{(\ell)}+R_a^{(\ell)}\right)}
                 {\sum_{\ell=1}^{m}\omega_\ell},
 \label{eq:twc-layer-objective}
\end{equation}
where $\omega_\ell$ is the weight assigned to layer $\ell$.
The last layer provides the final estimates, while earlier-layer
outputs provide intermediate supervision.
Gradients of \eqref{eq:twc-layer-objective} are propagated through
the denoisers, classifiers, and inter-layer messages. Thus activity
supervision also adjusts channel denoising, and channel losses at later
layers refine the activity estimates used in earlier feedback.

\subsection{Progressive Training}

Following the layer-wise training approach in \cite{wei2026}, we
increase receiver depth progressively. Each new layer has independent
network parameters initialized from the same pretrained score.
Its classifier correction and log calibration parameters are
initialized to zero, so the initial activity probability equals $\rho$.
Earlier layers retain their previously selected parameters.

At each depth, we first adapt the new layer and the feedback
calibration connecting it to the preceding layer. We then fine-tune all
layers jointly using a reduced learning rate. Both stages minimize
\eqref{eq:twc-layer-objective}. The checkpoint with the smallest
final-layer channel and activity loss on the validation set initializes
the next depth. Validation uses the training normalization constants.
